\section{Reference-line concentration and significant directions}
\label{conc:section}

We use the finite packet forms, their assignment to popular tops, and the
notation of the preceding sections.  In particular, all estimates concern
one of the separate forms in the packet expansion.  Its label multiplicity
and its spatial lattice multiplicity are bounded by absolute constants.
The Hilbert space of bin components is denoted by $\mathcal H$, and
$\norm{g(z)}_{\mathcal H}\leq\one_E(z)\leq1$.
Every subcollection in this section is fixed before its functions are
evaluated at a spatial point.

The first task is to control every interval of scales assigned to a top
outside one small exceptional set.  We prove moment bounds on
each reference line, then obtain the uniform top and derivative bounds
in \Cref{conc:top-good}.  \Cref{conc:mask-removal} uses these bounds to
replace the original adjoint sum $P$ by the masked sum $P_q$ and impose
the weighted top-count threshold.

For a slope $t$, put $Y_t(x,y)=y-tx$.  A $t$-box of dimensions $r\times W$
is a rectangle in the coordinates $(x,Y_t)$.  The affine change to these
coordinates preserves area.  Fix a reference length $l$ and consider
tiles satisfying
\begin{equation}\label{eq:src-9}
 \ell_s\leq l,\qquad
 \abs{\theta_s-t}\leq C_* B^4d_s,\qquad d_s=w/\ell_s.
\end{equation}
Here and below the fixed constant $C_*$ can be increased finitely many
times.  This only changes absolute constants in the conclusions.  The
inequality $B^{100}\ll m$ is more than sufficient for all such choices.

The tiles assigned to one top satisfy the narrower bound
$\abs{\theta_s-s_T}\leq C d_s$.  The wider $B^4$ window also covers the
fixed packet families used for residual estimates, where the reference
slope can be that of a different top; see
\eqref{hier:reference-packets} and \eqref{hier:deterministic-residual-subset}.

For a fixed, sufficiently large integer $n$, write
\begin{equation}\label{conc:envelope}
 \mathcal B_{s,n}^{t}(z)=
 \left(1+\frac{(x-x_s)^2}{\ell_s^2}\right)^{-n}
 \left(1+\frac{(Y_t(z)-Y_t(z_s))^2}{W^2}\right)^{-n}.
\end{equation}
The change from the packet coordinates of slope $\theta_s$ to these
coordinates has normalized shear at most $C_*B^4/m$.  Thus packet bounds
with any specified fixed number of derivatives and any specified fixed
decay order imply the corresponding bounds by
$C_n\mathcal B_{s,n}^{t}$.  We always choose the available packet decay
orders larger than the envelope orders used below.  These are fixed
integers; none grows with $L$, the number of scales, or $\dim\mathcal H$.

\subsection{Positive estimates and the main coefficients}

For a reference slope $t$, decompose the actual coefficient as
$a_s=a_s^{\rm main}+a_s^{\rm err}$, where componentwise
\begin{equation}\label{conc:main-coefficient}
 a_{s,j}^{\rm main}
 =\frac1{\abs{R_s}}\int
 \one_{\{\abs{u(z)-t}\leq d_s\}}
 g_j(z)\overline{b_{s,j}(t)\Psi_{s,j}(z)}\,\dd z.
\end{equation}
This convention agrees with the pairing, linear in its first argument,
used to define $a_s$.  In particular the error analyzes
$\overline{b_{s,j}(u)}g_j-
\one_{\{\abs{u-t}\leq d_s\}}\overline{b_{s,j}(t)}g_j$.

We first record a positive estimate which also applies when only one
scale is present.  For $r>0$ let
$k_{r,n}(x)=r^{-1}(1+x^2/r^2)^{-n}$.

\begin{lemma}\label{conc:single-scale}
For a fixed scale $\ell=w/d$, and any subcollection of tiles satisfying
\eqref{eq:src-9} at this scale,
\begin{equation}\label{conc:original-positive}
 \sum_{\ell_s=\ell}\norm{a_s}_{\mathcal H}
       \mathcal B_{s,n}^{t}(z)\leq C_n.
\end{equation}
There are fixed constants $0<c<1<C'$ for which
\begin{align}
 f_d(z)&=C\left(
 \one_{\{cd\leq\abs{u(z)-t}\leq C'B^4d\}}
 +\frac{\abs{u(z)-t}}{d}
       \one_{\{\abs{u(z)-t}\leq cd\}}\right),
       \label{conc:error-input}\\
 \sum_{\ell_s=\ell}\norm{a_s^{\rm err}}_{\mathcal H}
       \mathcal B_{s,n}^{t}(z)
 &\leq C_n
 \big(k_{\ell,n'}*_{x} k_{W,n'}*_{Y_t}f_d\big)(z)
       \label{conc:error-majorant}
\end{align}
for a fixed $n'>2$ with $2n'<n$.  The constants are independent of
the number of labels permitted in \eqref{eq:src-9}.
\end{lemma}

\begin{proof}
Minkowski's inequality bounds the norm of an analyzing coefficient by
the integral of its scalar analyzing envelope times the largest modulus
of its multiplier among the components, divided by $\ell W$; the factor
$\norm{g}_{\mathcal H}$ can be replaced by one.  At a fixed source point
the support condition for $b_{s,j}(u)$ allows only a bounded number of
labels at this scale.  This remains true after taking the maximum over
components, since the support condition is the common condition on the
label slope at precision $d$.

Choose the available analyzing decay order $N\geq2n$, with $n,n',N$
all fixed.  For one label, the lattice gives the unequal-exponent bound
\[
 \sum_{s\text{ in that lattice}}
 \frac{\mathcal B_{s,N}^{t}(z')\mathcal B_{s,n}^{t}(z)}{\ell W}
 \leq C_{n,N,n'} k_{\ell,n'}(x-x')
              k_{W,n'}(Y_t(z)-Y_t(z')).
\]
Indeed, in each normalized coordinate,
$1+\abs{a-b}^2\leq2(1+\abs{a-c}^2)(1+\abs{b-c}^2)$
extracts the relative kernel with exponent $n'$ from the two endpoint
envelopes.  The remaining exponents are $N-n'$ at the source and
$n-n'$ at the output.  Drop the source factor, which is at most one,
and sum the output envelope.  Its lattice sum is uniformly bounded
since $n-n'>1/2$: first sum the vertical centers in each horizontal row,
then the horizontal rows.  The transverse translations of the rows
do not affect these one-dimensional bounds.  The comparison with the
reference coordinates already gives the required analyzing envelope.
Summing the bounded number of labels active at the source and integrating
proves
\eqref{conc:original-positive}.

For the error, if $\abs{u-t}\leq cd$ with $c$ sufficiently small, the
replacement indicator is one and smoothness gives
$\abs{b_{s,j}(u)-b_{s,j}(t)}\leq C\abs{u-t}/d$.
The union of labels on which either multiplier is nonzero is bounded.
Outside this region, any nonzero original multiplier requires
$\abs{u-t}\leq C'B^4d$ by \eqref{eq:src-9}, whereas a nonzero replacement
requires $\abs{u-t}\leq d$.  The multipliers are bounded there.  Again,
only boundedly many labels are active at the source, counting both
families.  Applying the same lattice calculation with the common
factor \eqref{conc:error-input} proves \eqref{conc:error-majorant}.
\end{proof}

The smoothness and support assumptions on $b_{s,j}$ also show that
nonzero main coefficients use only a bounded number of labels per
scale.  In such a component,
$\abs{t-\theta_s}\asymp d_s$.  On a line of slope $t$, the Fourier
frequency of a packet is $\xi+t\eta$, and hence its support is contained
in
\begin{equation}\label{conc:main-annuli}
 c_a\ell_s^{-1}\leq\abs{\xi+t\eta}
                  \leq C_a\ell_s^{-1}
\end{equation}
with absolute $0<c_a<C_a$.  This uses the separation of the selector
from $\theta_s$ and the small choice of the packet bandwidth $c_*$.

Multiplying a synthesizing packet by its top mask preserves
\eqref{conc:main-annuli}, with slightly relaxed absolute constants.
Indeed the extra frequency on a $t$-line has magnitude at most
\begin{equation}\label{conc:mask-bandwidth}
 C L^{1-A}\left(l_{T(s)}^{-1}
             +\frac{\abs{t-s_{T(s)}}}{W}\right)
 \leq C L^{1-A}(1+B^4/m)\ell_s^{-1}.
\end{equation}
Here $\ell_s\leq l_{T(s)}$ and
$\abs{t-s_{T(s)}}\leq C B^4d_s$.  For $A\geq2$ and sufficiently large
$L$, this is less than $c_a/(10\ell_s)$.  Moreover, for every fixed
derivative order the masked packets have the same uniform envelope
bounds.  One can see this either from the explicit sinc product or
from its compact Fourier support and bounded supremum.  In particular
\begin{equation}\label{conc:mask-derivative}
 \abs{(\partial_x+t\partial_y)q_T}
 \leq C L^{1-A}\big(l_T^{-1}+\abs{t-s_T}/W\big),\qquad
 \abs{\partial_yq_T}\leq C L^{1-A}/W.
\end{equation}
These absolute bounds are global and do not require a good-point
assumption.

\begin{lemma}\label{conc:coefficient-carleson}
For every $t$-box $Q$ of width $W$, the main coefficients satisfy
\begin{equation}\label{eq:src-10}
 \sum_{\substack{z_s\in Q\\\ell_s\leq\mathrm{length}(Q)}}
       \norm{a_s^{\rm main}}_{\mathcal H}^2\abs{R_s}
       \leq C\abs{Q}.
\end{equation}
The same estimate holds after restricting to any fixed subcollection.
\end{lemma}

\begin{proof}
Let $\mathcal S_Q$ be the finite set in the sum and dualize its weighted
coefficient norm against vectors $c_s$ satisfying
$\sum_{s\in\mathcal S_Q}\norm{c_s}^2\abs{R_s}=1$.
Using \eqref{conc:main-coefficient}, the dual pairing is bounded by the
integral of the norm of the synthesis
\[
 \sum_{s\in\mathcal S_Q}
   \one_{\{\abs{u(z)-t}\leq d_s\}}
   \big(c_{s,j}b_{s,j}(t)\Psi_{s,j}(z)\big)_j.
\]
At each $z$, the indicators choose a partial interval of the dyadic
lengths: $\abs{u(z)-t}\leq w/\ell_s$ is monotone in $\ell_s$.
The norm of this synthesis is therefore at most its maximal partial
length sum.  All its nonzero component frequencies satisfy
\eqref{conc:main-annuli}.

Write $r=\mathrm{length}(Q)>0$.  Apply the proof of the Hilbert-valued
synthesis estimate in \Cref{tile:lem-synthesis}, including a fixed
polynomial growth weight $P_Q$ about $Q$.  This proof applies to the
present reference box even if $r$ is not dyadic or $r>\tau$.
For the weighted packet bounds it uses only center containment in $Q$
and $\ell_s\leq r$; multiplication by $P_Q$ preserves frequency support
and gives tile bounds with slower fixed decay.  The annular separation
and maximal partial sums depend on the dyadic tile lengths, not on the
reference length $r$.  For example, a fixed sufficiently large integer
$q$ gives the admissible polynomial
\[
 P_Q(z)=\left(1+\frac{(x-x_Q)^2}{r^2}\right)^q
        \left(1+\frac{(Y_t(z)-Y_t(z_Q))^2}{W^2}\right)^q,
 \qquad \int P_Q^{-2}\leq C_q rW=C_q\abs Q.
\]
The main label count is bounded, so the same synthesis proof yields
\[
 \int\norm{\text{selected synthesis}}
 \leq \left(\int P_Q^{-2}\right)^{1/2}
       \norm{P_Q\,\text{maximal synthesis}}_2
 \leq C\abs{Q}^{1/2}.
\]
Taking the supremum over $c_s$ proves \eqref{eq:src-10}.
Deleting packets preserves both the synthesis hypotheses and this
proof.
\end{proof}

\subsection{Fourier sums and a dyadic martingale}

For a finite sequence $(F_k)$ of functions with values in $\mathcal H$,
define its interval-sum variation by
\begin{equation}\label{conc:variation-definition}
 \mathcal V_3(F_k)(x)
 =\sup_{\mathcal I}
 \left(\sum_{I\in\mathcal I}
          \norm{\sum_{k\in I}F_k(x)}_{\mathcal H}^{3}\right)^{1/3},
\end{equation}
where $\mathcal I$ ranges over finite families of disjoint integer
intervals.  Empty intervals contribute zero.  This is bounded by the
usual $3$-variation of partial sums with the initial zero partial sum
included: insert the endpoints of the chosen intervals in order and
retain the resulting nonnegative terms.  This formulation allows the
interval endpoints to depend on the evaluation point.

\begin{lemma}\label{conc:annular-variation}
Suppose $F_k\in L^2(\R;\mathcal H)$ has Fourier support in
$\{\nu:c_a2^k\leq\abs\nu\leq C_a2^k\}$, with fixed positive
$c_a,C_a$.  Then
\begin{equation}\label{conc:annular-var-bound}
 \norm{\mathcal V_3(F_k)}_{L^2(\R)}
 \leq C_{c_a,C_a}
       \left(\sum_k\norm{F_k}_2^2\right)^{1/2}.
\end{equation}
The constant is independent of the finite index range and of
$\dim\mathcal H$.
\end{lemma}

\begin{proof}
We give the smooth-convolution versus dyadic-expectation comparison
explicitly; compare \cite[Lemma~3.2 and Section~4]{JSW2008}.
Split the integers into a fixed number $q$ of residue classes, large
enough that the annuli in one class have gaps of ratio greater than
four.  The variation of the sum is at most the sum of the variations
of these classes, so it suffices to treat one class.  Its supports are
disjoint, and $F=\sum_kF_k$ satisfies
$\norm F_2^2=\sum_k\norm{F_k}_2^2$.

Choose a smooth function $\chi$ equal to one on $[-1,1]$ and supported
in $[-2,2]$, and let $P_i$ have multiplier $\chi(2^{-i}\nu)$.
At a sequence of integer indices lying in the annular gaps, $P_iF$
is exactly the partial sum of all the preceding bands.  Include an
index below the first band, where this partial sum is zero.

For each integer $i$, let $\mathcal F_i$ be the sigma algebra generated
by the intervals
$[a2^{-i},(a+1)2^{-i})$, $a\in\Z$, and let $E_i$ be averaging on its
atoms.  The sigma algebras increase with $i$ and $(E_iF)_i$ is a
martingale.  We prove the comparison bound
\begin{equation}\label{conc:expectation-comparison}
 \sum_i\norm{(P_i-E_i)F}_2^2\leq C\norm F_2^2.
\end{equation}
Take a smooth dyadic partition $F=\sum_jf_j$ with $f_j$ supported in a
fixed annulus of size $2^j$ and
$\sum_j\norm{f_j}_2^2\leq C\norm F_2^2$.
If $j\leq i-C$, then $P_if_j=f_j$, and Poincar\'e's inequality on
each interval of length $2^{-i}$, followed by Bernstein's inequality,
gives
\[
 \norm{(P_i-E_i)f_j}_2\leq C2^{j-i}\norm{f_j}_2.
\]
If $j\geq i+C$, then $P_if_j=0$.  Define its band primitive $A_j$ by
$\widehat A_j(\nu)=\widehat f_j(\nu)/(i\nu)$.  Thus
$A_j'=f_j$ and $\norm{A_j}_2\leq C2^{-j}\norm{f_j}_2$.
With $h=2^{-i}$, averaging on a cell gives
$E_if_j=(A_j((a+1)h)-A_j(ah))/h$ there.
The sampling estimate
\begin{equation}\label{conc:band-sampling}
 \sum_{a\in\Z}\norm{A_j(ah)}_{\mathcal H}^{2}
 \leq C2^j\norm{A_j}_2^2,\qquad h\geq c2^{-j},
\end{equation}
follows directly from smooth reproduction.  Namely, take a Schwartz
kernel $K_j$ whose Fourier transform equals one on the spectrum of
$A_j$ and which satisfies
$\abs{K_j(x)}\leq C_N2^j(1+2^j\abs x)^{-N}$.
Cauchy--Schwarz in $A_j=K_j*A_j$ gives
$\norm{A_j(ah)}^2\leq C\int\abs{K_j(ah-x)}\norm{A_j(x)}^2\dd x$;
summing uses $\sup_x\sum_a\abs{K_j(ah-x)}\leq C2^j$.
Consequently
\[
 \norm{E_if_j}_2^2
 \leq C h^{-1}\sum_a\norm{A_j(ah)}^2
 \leq C2^{i-j}\norm{f_j}_2^2.
\]
For $\abs{i-j}\leq C$ use the bounded $L^2$ norms of $P_i,E_i$.
The resulting bounds have the form
$\norm{(P_i-E_i)f_j}_2\leq a_{i-j}\norm{f_j}_2$ with a fixed
$\ell^1(\Z)$ sequence $a$.  The triangle inequality followed by the
$\ell^1*\ell^2$ convolution inequality proves
\eqref{conc:expectation-comparison}.

For any sequence $(D_i)$, its usual $3$-variation is bounded pointwise
by $2(\sum_i\norm{D_i}^2)^{1/2}$: each endpoint of an increasing
sequence is used at most twice, and $\ell^3$ is bounded by $\ell^2$.
Thus \eqref{conc:expectation-comparison} controls the variation of
$(P_i-E_i)F$.

L\'epingle's martingale variation inequality \cite{Lepingle1976}, in
the unweighted $L^2$ case and at variation exponent $3$, gives
$\norm{V^3(E_iF)}_2\leq C\norm F_2$; see
\cite[Theorem~1.1, with $p=2$, $r=3$, and weight one]{ZorinKranich2020}.
For a finite index range on $\R$, apply its probability-space version
on every atom of the coarsest sigma algebra, normalizing Lebesgue
measure there, and sum the squared inequalities.  This proves the
same statement on the line without a finite-measure assumption.
For Hilbert-valued $F$, apply the scalar inequality to its real and
imaginary components.  Minkowski's inequality, in the form
$\ell^3(\ell^2)\leq\ell^2(\ell^3)$, and the outer $L^2$ norm give
the same constant independently of the number of components.
Combining with the comparison and then taking the selected low-pass
indices proves the estimate for one spaced class.  Summing the fixed
number of classes proves \eqref{conc:annular-var-bound}.
\end{proof}

\subsection{Every-line moment bounds}

We use the following precise local form of the John--Nirenberg
inequality \cite{JohnNirenberg1961}; a proof is included to specify
the normalization.

\begin{lemma}\label{conc:john-nirenberg}
Let $I$ be a finite interval and let $f$ be a locally integrable scalar
function on $I$.  Suppose that for each dyadic subinterval $J$ of $I$,
$\abs J^{-1}\int_J\abs{f-f_J}\leq A$, where $f_J$ denotes its average.
Then, for every $b\geq1$,
\[
 \left(\frac1{\abs I}\int_I\abs f^b\right)^{1/b}
 \leq\abs{f_I}+CbA.
\]
\end{lemma}

\begin{proof}
If $A=0$, dyadic differentiation makes $f$ constant almost everywhere,
and the assertion is immediate.  Otherwise take the maximal proper
dyadic subintervals $J$ of $I$ on which the average of
$\abs{f-f_I}$ exceeds $2A$.  Their total length is at most
$\abs I/2$, the function has $\abs{f-f_I}\leq2A$ almost everywhere
off their union, and dyadic parent maximality gives
$\abs{f_J-f_I}\leq4A$ on each selected interval.  Repeat the same
selection inside each $J$, now subtracting its own mean.  At generation
$k$ the selected union has measure at most $2^{-k}\abs I$.
Outside it, the successive differences of means and the final
differentiation bound show
$\abs{f-f_I}\leq(4k+2)A$.
This proves an exponential distributional bound for $\abs{f-f_I}/A$.
Integrating that bound by layer cake gives
$\norm{f-f_I}_{L^b(I,\dd x/\abs I)}\leq CbA$; adding the mean proves
the lemma.  Finite stopping generations and then monotone convergence
justify the argument for any locally integrable $f$.
\end{proof}

Fix a sufficiently large envelope exponent $n$.  At each scale put
\[
 U_\ell(z)=\sum_{\ell_s=\ell}\norm{a_s^{\rm main}}\mathcal B_{s,n}^{t}(z),
 \qquad
 \mathcal E_t(z)=\sum_s\norm{a_s^{\rm err}}\mathcal B_{s,n}^{t}(z),
 \qquad
 \mathcal S_t(z)=\left(\sum_\ell U_\ell(z)^2\right)^{1/2}.
\]
These quantities can be formed on any fixed subcollection; it is also
permissible to dominate a subcollection by the positive quantities on
the full collection.

\begin{proposition}\label{conc:line-moments}
For every slope $t$, every collection satisfying \eqref{eq:src-9},
every line $Y_t=y_0$, every horizontal interval $I$ of length $l$ on
that line, and every $b\geq2$, the normalized $L^b(I)$ norms of
$\mathcal E_t$, $\mathcal S_t$, and
\[
 \mathcal V_3\left(\sum_{\ell_s=2^k}\phi_s\right),
 \qquad
 \mathcal V_3\left(\sum_{\ell_s=2^k}q_{T(s)}\phi_s\right)
\]
are at most $C b^C(1+\log B)^C$.
The constants are uniform in the line, interval, finite collection,
dimension, and number of scales.  In particular the assertion holds
for every fixed deterministic subcollection with the same constants.
\end{proposition}

\begin{proof}
We give the local mean and oscillation bounds, including their behavior
on intervals shorter than $l$.  All convolutions in what follows are
in the $(x,Y_t)$ coordinates.

First, \eqref{conc:error-input} satisfies
\begin{equation}\label{conc:error-input-sum}
 0\leq f_d\leq C,\qquad \sum_{d\text{ dyadic}}f_d\leq C(1+\log B).
\end{equation}
The first indicators are nonzero for only $O(1+\log B)$ dyadic $d$ at
a fixed point, and the second terms sum geometrically.  Average first
with the common vertical kernel $k_{W,n'}$.  Its $L^1$ norm is fixed,
so \eqref{conc:error-input-sum} persists on each reference line.
The right side of \eqref{conc:error-majorant}, summed over scales,
therefore defines a nonnegative majorant $\mathcal A_t$ for
$\mathcal E_t$.

Let $J$ be any horizontal interval of length $r\leq l$.  If $\ell\leq r$,
\[
 \frac1r\int_J k_{\ell,n'}(x-x')\dd x
 \leq \frac Cr\left(1+\frac{\dist(x',J)}r\right)^{-2}.
\]
Hence the local mean of the short-scale part of $\mathcal A_t$ is at
most $C(1+\log B)$, by summing the source functions before integration.
For $\ell>r$, the derivative of its convolution with a bounded source
has supremum at most $C/\ell$.  The total oscillation of these terms
on $J$ is at most $Cr\sum_{\ell>r}\ell^{-1}\leq C$.
It follows that $\mathcal A_t$ has local BMO norm
$O(1+\log B)$.  Its mean on an interval of length $l$ has the same
bound, since all scales then belong to the short-scale part.
John--Nirenberg on that interval gives
$\norm{\mathcal E_t}_{L^b(I,\dd x/l)}\leq Cb(1+\log B)$.
We only use domination by the BMO majorant here; no assertion that an
arbitrary function below a BMO function is BMO is required.

For the square function, fixed-scale lattice overlap gives
\begin{equation}\label{conc:lattice-square}
 U_\ell(z)^2
 \leq C\sum_{\ell_s=\ell}
             \norm{a_s^{\rm main}}^2\mathcal B_{s,n}^{t}(z).
\end{equation}
For the same interval $J$, partition center space into $t$-boxes
$Q_{ab}$ of dimensions $r\times W$, with horizontal offset $a$ from
$J$ and transverse offset $b$ from the reference line.  For
$z_s\in Q_{ab}$ and $\ell_s\leq r$, the envelope bound implies
\[
 \int_J\mathcal B_{s,n}^{t}(x,y_0+tx)\dd x
 \leq C\ell_s(1+\abs a+\abs b)^{-n_1}
\]
for a fixed $n_1>4$, after choosing $n$ sufficiently large.
Indeed the transverse decay gives the $b$ factor; for a distant
horizontal box the integral is bounded by
$Cr(\ell_s/(r\abs a))^{n_1}$, which is at most
$C\ell_s(1+\abs a)^{-n_1}$ because $\ell_s\leq r$.
Equation~\eqref{eq:src-10} gives exactly
\[
 \sum_{\substack{z_s\in Q_{ab}\\\ell_s\leq r}}
       \norm{a_s^{\rm main}}^2\ell_s
 =W^{-1}\sum_{\substack{z_s\in Q_{ab}\\\ell_s\leq r}}
       \norm{a_s^{\rm main}}^2\abs{R_s}\leq Cr.
\]
Summing the offsets in \eqref{conc:lattice-square} proves
\begin{equation}\label{conc:line-short-square}
 \frac1r\int_J\sum_{\ell\leq r}U_\ell^2\dd x\leq C.
\end{equation}
This argument applies on each line directly; it is not a restriction
of an almost-everywhere plane estimate to a line.  Notice also that
the factor $W$ has canceled, so the estimate has no dependence on $m$.
Since main coefficients are bounded, fixed-scale lattice bounds give
$U_\ell\leq C$ and $\abs{\partial_x U_\ell}\leq C/\ell$ on the line.
The long-scale part of $\mathcal S_t^2$ has oscillation at most $C$ on
$J$.  By \eqref{conc:line-short-square}, its local BMO norm is bounded;
its mean on every length-$l$ interval is bounded as well.
John--Nirenberg applied to $\mathcal S_t^2$ gives
$\norm{\mathcal S_t}_{L^b(I,\dd x/l)}\leq C\sqrt b$.

It remains to prove the main-term variation bound, with or without
the mask.  Fix $J$ and its center boxes $Q_{ab}$ as above, and first
restrict to scales $\ell_s\leq r$.  Group the synthesis by those boxes.
For a fixed box multiply its synthesis by
\[
 P_{ab}(x)=\left(1+\frac{(x-x_{Q_{ab}})^2}{r^2}\right)^{q_0},
\]
where $q_0$ is a sufficiently large fixed integer.  Because
centers lie in $Q_{ab}$ and $\ell_s\leq r$,
these multiplied waves still have fixed tile envelope bounds.
Their Fourier supports are preserved by polynomial multiplication.
Using \eqref{conc:lattice-square} and integrating on the entire line
therefore gives
\[
 \sum_\ell\left\|P_{ab}
   \sum_{\substack{\ell_s=\ell\\z_s\in Q_{ab}}}
       a_s^{\rm main}\Phi_s\right\|_{L^2(\R;\mathcal H)}^2
 \leq Cr(1+\abs b)^{-n_1},
\]
and the same bound holds with $q_{T(s)}\Phi_s$ in place of $\Phi_s$.
Here and elsewhere $a_s^{\rm main}\Phi_s$ denotes componentwise
multiplication.  The bound follows from \eqref{eq:src-10} as above,
retaining transverse decay while integrating the horizontal envelope.
Apply \Cref{conc:annular-variation} with the length index reversed;
this preserves the interval-sum variation.  The support hypotheses hold by
\eqref{conc:main-annuli} and \eqref{conc:mask-bandwidth}.  On $J$,
$P_{ab}\geq c(1+\abs a)^{2q_0}$, and multiplication by this one scalar
commutes with every interval sum and with its norm.  Thus the
$L^2(J)$ variation of the synthesis from $Q_{ab}$ is at most
\[
 Cr^{1/2}(1+\abs a)^{-2q_0}(1+\abs b)^{-n_1/2}.
\]
Minkowski and summation over $a,b$ give a bound $Cr^{1/2}$ for the
short-scale main variation.

For each longer scale the derivative of main synthesis on the line
is bounded by $C/\ell$, also after masking by
\eqref{conc:mask-derivative}.  The variation of the long-scale part
therefore changes by at most
$Cr\sum_{\ell>r}\ell^{-1}\leq C$ on $J$.
The full main variation differs from the long-scale variation by at
most the short-scale variation pointwise, by the triangle inequality
for \eqref{conc:variation-definition}.  This proves its bounded local
BMO norm.  Taking $r=l$ proves its bounded mean on $I$.
John--Nirenberg gives normalized $L^b$ norm at most $Cb$.
The variation of the error synthesis is bounded by the positive sum
$C\mathcal E_t$, since disjoint intervals use each scale at most once.
For the masked error the same is true because $\abs{q_T}\leq1$.
The already proved error bound completes the proposition.
\end{proof}

All arguments in \Cref{conc:line-moments} use only fixed support,
decay, and derivative bounds on the synthesizing waves.  They still
hold if those waves are replaced by waves with the same supports and
such bounds multiplied by a fixed constant $C_0$; the conclusions are
then multiplied by $C C_0$.  This observation includes fixed polynomial
growth weights about a region containing the centers, provided the
polynomial is absorbed into a slower fixed tile envelope.  It also
includes the componentwise operator $W(\partial_y-i\eta_j)$ on the
unmasked synthesizing waves.  After demodulation its multiplier has
bounded magnitude on the vertical packet support.  These observations
will be used only at fixed orders below.

\subsection{Uniform top bounds and neighborhood derivatives}

Choose a fixed exponent $A_1$ so that $K_1=L^{A_1}$ contains the
polynomial dilation of $R_T$ in which all assigned centers lie, with a
fixed additional margin.  This choice depends on the top assignment
and the fixed weight exponent $D_2$, and is made before choosing $A$.
Set
\[
 h_T^{(1)}(z)=C_3\big(1+\dist_T(z)/K_1\big)^{-D_2}.
\]
For an interval $\mathcal J$ of dyadic lengths define the top sum
$F_{T,\mathcal J}=\sum_{s:T(s)=T,\ \ell_s\in\mathcal J}\phi_s$.
The family of these sums includes empty intervals and is finite for
our finite packet collection.  Let
$\mathcal D F=(W(\partial_y-i\eta_j)F_j)_j$.

\begin{corollary}\label{conc:top-good}
For every fixed $J_*>0$ there are a fixed exponent $C_{\mathrm{top}}$ and a
measurable set $\mathcal X_{\rm top}$ with
$\abs{\mathcal X_{\rm top}}\leq S e^{-J_*L}$ such that, for all tops
$T$ and all $z\notin\mathcal X_{\rm top}$,
\begin{equation}\label{eq:src-11}
 \sup_{\mathcal J}\ \sup_{\abs v\leq W}
 \left(\norm{F_{T,\mathcal J}(z+(0,v))}_{\mathcal H}
       +\norm{\mathcal D F_{T,\mathcal J}(z+(0,v))}_{\mathcal H}\right)
 \leq L^{C_{\mathrm{top}}}\big(h_T^{(1)}(z)\big)^8.
\end{equation}
Independently of this exceptional set, for every $z$, every $l_*>0$,
and every interval $\mathcal J\subset[l_*,\infty)$,
\begin{equation}\label{conc:top-horizontal-derivative}
 \norm{(\partial_x+s_T\partial_y)F_{T,\mathcal J}(z)}_{\mathcal H}
 \leq \frac{L^{C_{\mathrm{top}}}}{l_*}\big(h_T^{(1)}(z)\big)^8.
\end{equation}
All packet orders, $K_1$, and $C_{\mathrm{top}}$ can be fixed before $A$;
they are independent of the later mass cutoff and hierarchy depth.
\end{corollary}

\begin{proof}
Use \Cref{conc:line-moments} on assigned top tiles, with
$t=s_T$ and $l=l_T$.  These satisfy \eqref{eq:src-9} with a fixed
angular constant.  Write $X=(x-x_T)/l_T$ and
$Y=(Y_{s_T}(z)-Y_{s_T}(z_T))/W$.  For a sufficiently large fixed integer
$q$, let
\[
 P_T(z)=\left(1+\frac{X^2+Y^2}{K_1^2}\right)^q.
\]
Its product with any assigned wave is bounded by a slower fixed tile
envelope with uniform constants.  Indeed at the centers $X,Y$ are
bounded by a fixed fraction of $K_1$, and in moving away from a center
the top scales are no smaller than the corresponding tile scales.
The same statement holds for the fixed derivatives in use.
Multiplication by $P_T$ preserves spectrum.  Applying the last
observation after \Cref{conc:line-moments} gives a uniform polynomial
moment bound on every length-$l_T$ interval of every parallel line
for
$\sup_{\mathcal J}\norm{P_TF_{T,\mathcal J}}$ and
$\sup_{\mathcal J}\norm{P_T\mathcal D F_{T,\mathcal J}}$.

We justify the required vertical neighborhood supremum.  Demodulate
component $j$ by $e^{-i\eta_jY_{s_T}}$.  Each of the just described
weighted sums has vertical Fourier support in $[-C/W,C/W]$.
Choose one Schwartz reproduction kernel $K_W$ with Fourier transform
one on that interval.  Its absolute values at all arguments shifted
by at most $W$ are bounded by
$C_nW^{-1}(1+\abs y/W)^{-n}$.  Reproduction, the Hilbert norm
inequality, and the supremum over $\mathcal J$ show that its vertical
neighborhood supremum is bounded by convolution of the corresponding
unshifted supremum with this positive integrable kernel.  Minkowski
in normalized $L^b$ on horizontal intervals preserves the uniform
moment estimate.  Finally $P_T(z+(0,v))\asymp P_T(z)$ for
$\abs v\leq W$, uniformly in $z$, because $K_1\geq1$.
We have therefore proved, on every top-scale box, polynomial moments
for $P_T$ times the entire left side of \eqref{eq:src-11}.

For clarity, this yields a global exceptional-area bound, not merely
one inside the top box.  Partition the plane into boxes of dimensions
$l_T\times W$, indexed by $(a,b)\in\Z^2$.  On each such box,
$1+\dist_T/K_1\asymp1+(\abs a+\abs b)/K_1$.
Choose $q$ with $2q>8D_2+4$.  If the left side of
\eqref{eq:src-11} exceeds its right side in that box, its product
with $P_T$ exceeds
\[
 c L^{C_{\mathrm{top}}}
       \big(1+(\abs a+\abs b)/K_1\big)^{2q-8D_2}.
\]
Take moment order $b_L=\lceil L\rceil$.  The available moment bound is
at most a fixed power of $L$, since $1+\log B$ is such a power.
Choose $C_{\mathrm{top}}$ larger by a fixed amount.  Markov's inequality
and summation over $(a,b)$ then give
\[
 \abs{\mathcal X_T}
 \leq C\abs{R_T}K_1^2 e^{-(J_*+1)L}
 \leq \abs{R_T}e^{-J_*L}
\]
for sufficiently large $L$, after increasing the threshold exponent
or the fixed lower bound on $L$.  Equivalently one can first arrange
$e^{-(J_*+2)L}$ to absorb the polynomial $K_1^2$.
Taking the union over tops proves the asserted bound by
$S=\sum_T\abs{R_T}$.  This argument uses fixed $q$ and fixed packet
orders, even though the moment order grows with $L$.

For \eqref{conc:top-horizontal-derivative}, coefficient norms are
bounded, there are boundedly many assigned labels per scale, and
each differentiated wave has a bound $C\ell_s^{-1}$ times a fixed
tile envelope in slope $s_T$.  At a fixed scale, summing this envelope
over assigned centers is at most
$C\ell_s^{-1}(1+\dist_T/K_1)^{-8D_2}$.
To see the decay, multiply by the inverse top factor, absorb it into
a slower tile envelope as for $P_T$, and use bounded lattice overlap.
Summing over $\ell_s\geq l_*$ costs
$\sum_{\ell\geq l_*}\ell^{-1}\leq C/l_*$.
The same bound applies to any interval of these lengths by the
triangle inequality.  This proves the unconditional derivative
estimate and completes the corollary.
\end{proof}

\subsection{The masked sum and the count threshold}

The choice of $A$ in the preceding section can now be completed.
Its dependence on the fixed constants just obtained is allowed;
their construction did not depend on $A$.

\begin{proposition}\label{conc:mask-removal}
Given any fixed $H>0$, choose $A$ sufficiently large after $K_1$ and
$C_{\mathrm{top}}$, and then choose a sufficiently large fixed count
threshold exponent $C_n$.  There is a measurable set
$F_{\rm conc}\subset F$ such that
\[
 \abs{F\setminus F_{\rm conc}}\leq L^{-H}\abs F,
 \qquad
 \int_{F_{\rm conc}}\norm{P-P_q}_{\mathcal H}
       \leq L^{-H}N\abs F,
 \qquad P_q=\sum_sq_{T(s)}\phi_s,
\]
and \eqref{eq:src-11} holds there for every top.  On this set,
\begin{equation}\label{eq:src-12}
 n(z)=\sum_T h_T(z)\leq NL^{C_n}.
\end{equation}
One has $K_0=L^{A+2}>K_1$, so both
\eqref{eq:src-11} and \eqref{conc:top-horizontal-derivative} also hold
with $h_T$ in place of $h_T^{(1)}$.
\end{proposition}

\begin{proof}
First fix, for example, $J_*=10$ in \Cref{conc:top-good}, and fix
the resulting $C_{\mathrm{top}}$ before making any choice of $A$.
Its exceptional area is smaller than any prescribed fixed inverse
power of $L$ times $\abs F$, since
$S/\abs F\leq L^C N^2\leq L^C e^{4L}$.
On $\dist_T\leq L^{A/2}$, the elementary estimate for the mask is
$\abs{1-q_T}\leq C L^{1-A}$: use
$\abs{1-\operatorname{sinc}v}\leq Cv^2$ near zero and the integer
power $2\lceil L\rceil$ in \eqref{eq:src-8}.
Outside this region, $\abs{1-q_T}\leq2$, whereas
\[
 \big(h_T^{(1)}\big)^7
 \leq C L^{-7D_2(A/2-A_1)}.
\]
Equation~\eqref{eq:src-11}, with the full length interval, therefore
implies outside $\mathcal X_{\rm top}$ that, for any prescribed fixed
$H_1$,
\begin{equation}\label{conc:mask-pointwise}
 \norm{P-P_q}\leq L^{-H_1}\sum_T h_T^{(1)}
\end{equation}
once $A$ is chosen large enough.  In the near region the factor
$L^{C_{\mathrm{top}}+1-A}$ suffices; in the far region use the displayed
seventh power to reduce $(h_T^{(1)})^8$ to $h_T^{(1)}$.

Apply \eqref{eq:src-6} with dilation $K_1$ to
$n_1=\sum_Th_T^{(1)}$.  For any finite-measure set $A_0$ in the testing
domain, weak $L^2$ and layer cake imply
\[
 \int_{A_0}n_1\leq2\norm{n_1}_{2,\infty}\abs{A_0}^{1/2}
 \leq L^C S^{1/2}\abs{A_0}^{1/2}.
\]
Since $S\leq L^C\abs E$ and $\abs E=N^2\abs F$, choosing $H_1$
larger than $H$ by the fixed powers in these inequalities proves the
claimed mask integral even when $A_0=F\setminus\mathcal X_{\rm top}$.

After $A$ has been fixed, \eqref{eq:src-6} with dilation $K_0$ gives
\[
 \abs{\{z\in F:n(z)>NL^{C_n}\}}
 \leq L^{C-2C_n} S/N^2
 \leq L^{C-2C_n}\abs F.
\]
Take $C_n$ sufficiently large and define $F_{\rm conc}$ by deleting
this set and $\mathcal X_{\rm top}$ from $F$.  Increasing the fixed
thresholds slightly makes the sum of their areas at most
$L^{-H}\abs F$.  This proves all assertions.
\end{proof}

The original adjoint estimate \eqref{eq:src-4} gives
$\norm P_2\leq L^{C_P}\abs E^{1/2}$ for a fixed exponent $C_P$.
Thus, by taking $H$ as large as required, any discarded set above has
arbitrarily small prescribed polynomial pairing with the original
$P$.  The concentration estimates above hold for finite tile sets and
scale ranges, with constants independent of their cardinalities.  After
the later assembly proves the uniform model bound for each finite form,
the estimate passes to the convergent bilinear packet representation as
in \Cref{red:packet-representation,red:model-suffices}.
